\documentclass[12pt]{article}
\usepackage{amsmath, amssymb, amsthm}
\usepackage{parskip}
\usepackage{geometry}
\usepackage{enumitem}
\usepackage{fancyhdr}

% Set margins to match a standard homework document
\geometry{a4paper, margin=1in}

\title{Recitation 10}
\parskip=1ex
\author{Math 311 Fall 2026 \\ Rutgers University}
\date{}

\begin{document}

\maketitle

\textbf{Exercise 1}
\begin{enumerate}[label=(\alph*)]
    \item A set \( O \) is open iff for every \( x \in O \), there exists \( V_{\epsilon}(x) \subseteq O \).

          A point \( x \) is a limit point of \( A \) iff for every \( \epsilon > 0 \) we have
          \[
              (V_{\epsilon}(x) \setminus {x}) \cap A \neq \varnothing .
          \]

          A point \( x \) is an isolated point iff there exists \(\epsilon > 0\) such that
          \[
              V_{\epsilon}(x) \cap A = {x}.
          \]

          A set \( A\) is closed if it contains all of its limit points.

          The closure subset of a set \( A \) is defined as \( \overline{A} = A \cup L \) where \( L \) is the set of all of \(A\) limit points.

          No, the limit points of a set does not need to belong to a set.

    \item A set \( A \subseteq \mathbb{R} \) is a compact set iff every sequence \( a_n \in A \) has a subsequence \(x_{n_k}\) that converges in \(A\).

          An open cover is a collection of open sets \( O_\lambda : \lambda \in \Lambda \) such that
          \[
              A \subseteq \cup_{\lambda \in \Lambda} O_\lambda.
          \]

          The finite sub cover is any finite subset of the collection of open sets that make up the open cover such that the sub cover still contains \( A \).

    \item A set \( A \) is perfect iff it is closed and has no isolated points

          Sets \(A, B \) are separated iff \( \overline{A} \cap B = \varnothing \) and \(A \cap \overline{B} = \varnothing \).

          A set \( E \) is disconnected iff \(E = A \cup B\) and \(A,B\) are separated and nonempty.

          A set \( E \) is connected iff it is not disconnected

    \item A set \( A \) is dense iff \( \overline{A} = \mathbb{R} \).

          A set \( A \) is nowhere dense iff \(\overline{A}\) has no nonempty open intervals.

          A set \( A \) is \( F_\sigma \) iff it can be written as the infinite union of closed sets.

          A set \( A \) is \( G_\delta \) iff it can be written as the infinite intersection of open sets.

    \item A set \( O\) is open iff its complement \( O^c \) is closed and a set \( A \) is closed iff its complement \( A^c\) is open.

          Proof:
          For the forward direction, assume \( O \) is an open set. Let \( x \) be a limit point of \( O^c \). Then every neighborhood of \(x\) must intersect with \(O^c\).
          Suppose for contradiction, that \(x \in O\). Since \( O\) is an open set, there exists \(V_\epsilon(x) \subseteq O\). This contradicts the fact that every neighborhood of \(x\) must intersect with \(O^c\). Hence \(x \in O^c\). Since \(x\) is an arbitrary limit point, it follows that \( O^c \) is closed.

          Conversely, assume \( O^c \) is a closed set. Let \(x \in O\). Then \(x\) cannot be a limit point of \(O^c\). So there is exists a neighborhood of \(x\) that does not intersects with \(O^c\). Hence \( V_\epsilon(x) \subseteq O \). Since \( x \in O\) was arbitrary, it follows that \( O \) is open.

          Then for the second statement we can just take the complement of \( O \) for both directions.

          A set \(A\) is closed iff every convergent sequence \( a_n \in A \) converges to a point in \( A \).

    \item We know that \( \overline{A} \) is closed and contains all the limit points of \(A\) and that \( A \subseteq \overline{A} \). We know that \( \overline{A} \) is the smallest closed set containing \(A\).

    \item Heine-Borel Theorem states that the following three statements are equivalent.

          1. The set \( A \) is compact

          2. The set \( A \) is closed and bounded

          3. Every open cover of \( A \) has a finite subcover.

    \item The intersection
          \[
              \bigcap_{n=1}^{\infty} K_n \neq \varnothing.
          \]

    \item The set \( P \) is uncountable.

    \item A set \( E \) is connected iff for every \(a, b \in E\) such that \( a < b\), every \( c \in \mathbb{R} \) where  \( a < c < b\), we have \(c \in E\).

          Suppose \( E = A \cup B\) where \(A, B\) are nonempty and disjoint sets. Then \( E \) is connected if there exists a sequence in either \( A \) or \( B\) that converges to a point in the other set.

    \item The intersection of countably many dense open subsets of \( \mathbb{R} \) is nonempty.

          Baire's Theorem states that the reals \( \mathbb{R} \) cannot be written as a countable union of nowhere-dense sets.

    \item Given \(f: A \to \mathbb{R} \) and limit point \(c\) of \(A\).
          A functional limit \( \lim_{x \to c} f(x) = L \) exists iff for all sequences \( x_n \in A \) where \(x_n \neq c\) and \(x_n \to c \) we have \( f(x_n) \to L\).

          Proof: Suppose \( \lim_{x \to c} f(x) = L \). Let \( x_n \in A \) with \(x_n \neq c\) and \(x_n \to c \), and let \( \epsilon > 0 \). Since \( \lim_{x \to c} f(x) = L \), it follows that for all \(x\) satisfying \( 0 < | x - c | < \delta \), we have \( | f(x) - L | < \epsilon \). Choose \(N \in \mathbb{N} \) big enough such that for all \( n > N \), we have
          \[
              0 < | x_n - c| < \delta.
          \]
          Then we also have
          \[
              | f(x_n) - L | < \epsilon.
          \]
          That is, \( f(x_n) \to L\).

          Conversely, suppose  \( f(x_n) \to L\). Then suppose, for contradiction, that \( lim_{x \to c} f(x) \neq L \). So there exists \( \epsilon_0 > 0\) such that for every \( \delta > 0 \), there exists \(x \in A\) such that \( 0 < | x - c| < \delta \) but \( |f(x) - L | \geq \epsilon_0 \). For each \( n \in \mathbb{N} \), let
          \(\delta_n = 1/n\). Then there exists \(x_n \in A\) such that \( 0 < | x_n - c| < \delta_n \) and \( |f(x_n) - L | \geq \epsilon_0 \). Hence \(f(x_n) \not \to L\), contradicting our original assumption. Therefore, \( lim_{x \to c} f(x) = L \).

          A functional limit \( \lim_{x \to c} f(x) = L \) does not exist if there exists two sequence \( x_n, y_n \in A\) such that \(x_n,y_n \to c\) but
          \[
              \lim_{n \to \infty} f(x_n) \neq \lim_{n \to \infty} f(y_n).
          \]

    \item Given \(f: D \to \mathbb{R} \) and \( c\in D\). The function \(f\) is continuous at point \(c\) iff for every \(\epsilon > 0\), there exists \(\delta >0\) such that if \( | x - c | < \delta \), then \( |f(x) - f(c) | < \epsilon \). Moreover, the function \(f\) is continuous on its domain, if it is continuous on every point in its domain \(D\).

          The sequential definition states that \(f\) is continuous as point \(c\) iff for every sequence \( x_n \in A\) with \( x_n \to c \), we have \( f(x _n) \to f(c) \).

          If \(c\) is an isolated point, then it is always continuous since \( x_n = c \) will satisfy the sequential definition.

    \item Suppose \( f,g : d \to \mathbb{R} \) is continuous at \( c \in D\).
          The scalar-multiple rule says that \(k f(x) \) is continuous at \(c\) for all \( k \in \mathbb{R} \). The sum rule says that \( f(x) + g(x) \) is continuous at \(c\). The product rule says that \(f(x)g(x)\) is continuous at \(c\). The quotient rule says that
          \(\frac{f(x)}{g(x)}\) is continuous at \(c\) as long as \(g(c) \neq 0\).

          Let \(f: D \to \mathbb{R} \) and \(h: E \to \mathbb{R} \) satisfy \(f(D) \subseteq E\).
          If \(f\) is continuous at \(c\) and \(h\) is continuous at \(f(c\) , then the composition \(h \circ f\) is continuous at \(c\).

    \item Let \(f: D \to \mathbb{R} \). Then \(f\) is uniformly continuously iff for every \(\epsilon > 0\), there exists \(\delta > 0\) such that for every \(x,y \in D\)
          \[
              | x - y | < \delta \implies |f(x) -f(y) | < \epsilon.
          \]

    \item Suppose \(f: D \to \mathbb{R} \) is continuous and \(K \subseteq D\) is compact. Then \(f(K)\) is also compact.

          The complete theorem states that if \(f: D \to \mathbb{R} \) is a continuous function and \(K \subseteq D\) is compact, then \(f(K)\) is also a compact set.

    \item Suppose \(K\subseteq\mathbb R\) is nonempty and compact, and \(f:K\to\mathbb R\) is continuous. By Heine-Borel, \(f(K)\) is closed and bounded. It's supremum and infinum are both attained. By definition of image, there exists \(x_{max}, x_{min} \in K\) where \(f(x_{max}) = sup(f(K) \) and \(f(x_{min}) = inf(f(K))\).

    \item Suppose \(K\subseteq\mathbb R\) is compact and \(f:K\to\mathbb R\) is continuous. Then \(f\) is uniformly continuous on \(K\).

    \item Suppose \(D\subseteq\mathbb R\) is bounded and \(f:D\to\mathbb R\) is uniformly continuous. Then \(f(D)\) is also bounded.

    \item Suppose \(K\subseteq\mathbb R\) is compact and
          \[
              f:K\to\mathbb R
          \]
          is continuous and one-to-one. Then the inverse \(f^{-1}: f(K) \to K\) exists and is continuous.
\end{enumerate}

\textbf{Exercise 2}
\begin{enumerate}[label=(\alph*)]
    \item False. Consider \(f(x) = x^2\) and \(U = (-1,1)\). Then \(f(U) = [0, 1)\) which is not open.

    \item False. Consider
          \[
              f(x) = \frac{1}{1+x^2}
          \]
          and \(F = [0, \infty)\). Then \(f(F) = (0,1] \) which is not closed.

    \item True. Let \(y \in f(\overline{A})\). Then there exists \(x \in \overline{A} \) such that \(f(x) = y\). Moreover, there exists a sequence \(x_n \in A \) where \(x_n \to x\). Since \(f\) is continuous, it follows that \(f(x_n) \to f(x) = y\). Since \(x_n \in A\), we also have \(f(x_n) \in f(A)\) for every \(n \in \mathbb{N}\). Thus \(y \in \overline{f(A)}\). Therefore
          \[
              f(\overline{A}) \subseteq \overline{f(A)}.
          \]

    \item False. Consider
          \[
              f(x)=
              \begin{cases}
                  1 & x=0,    \\
                  0 & x\neq0.
              \end{cases}
          \]
          Then for any compact set \(K \subseteq \mathbb{R}\),
          \[
              f(K) \subseteq \{0,1\}
          \]
          Since \(f(K)\) is finite, it must be compact. Now consider the sequence \(x_n = 1/n\).
          Notice \(x_n \to 0\). But \(f(x_n) = 0 \not \to 1 = f(0)\). Thus \(f\) is not continous at 0.

    \item True. Suppose \(f,g:\mathbb R\to\mathbb R\) are continuous and
          \[
              f(r)=g(r)\qquad\forall r\in\mathbb Q.
          \]
          Let \(x \in \mathbb{R}\). Since \(\mathbb{Q} \) is dense in \(\mathbb{R}\), there exists a sequence \(r_n \subseteq \mathbb{Q}\) such that \(r_n \to x\). By continuity of \(f\) and \(g\),
          \[
              f(r_n) \to f(x), \qquad g(r_n) \to g(x).
          \]
          From the original assumption we have
          \[
              f(r_n) = g(r_n)
          \]
          for all \(n \in \mathbb{N}\).
          Therefore \(f(x) = g(x)\) for all \(x \in \mathbb{R}\).

    \item False. Let
          \[
              f(x) = g(x) = x
          \]
          and \(D = (0,1) \). Then \(T = (0,1) \) which is not closed.

    \item False. Notice that \(A\) is the union of two connected intervals, \((0,1)\) and \((1,2)\). Since \(f\) is continuous, the images \(f((0,1))\) and \(f((1,2))\) are each connected and hence intervals.
          Thus,
          \[
              f(A)=f((0,1))\cup f((1,2))
          \]
          is the union of at most two intervals. However,
          \[
              (-1,1)\cup(2,3)\cup(4,5)
          \]
          has three connected components and cannot be written as the union of two intervals. Therefore, no such continuous function exists.
\end{enumerate}

\textbf{Exercise 3}
\begin{enumerate}[label=(\alph*)]
    \item Suppose \(f:\mathbb R\to\mathbb R\) is continuous and \(U\subseteq\mathbb R\) is open. Let \(x \in f^{-1}(U)\). By definition of preimage, \(f(x) \in U\). Since \(U\) is open, there exists \(\epsilon > 0\) such that
          \[
              V_{\epsilon}(f(X)) \subseteq U.
          \]
          Since \(f\) is continuous, there exists \(\delta > 0\) such that
          \[
              | y -x | < \delta \implies | f(y) - f(x) | < \epsilon.
          \]
          Suppose \(y \in V_{\delta}(x)\). Then \(| f(y) - f(x) | < \epsilon\). Hence \(f(y) \in V_{\epsilon}(f(x))\). By definition of preimage, \(y \in f^{-1}(U)\). Thus
          \[
              V_{\delta}(x) \subseteq f^{-1}(U).
          \]
          Since \(x\) was arbitrary, we conclude that \(f^{-1}(U)\) is open.

    \item  Suppose \(f:\mathbb R\to\mathbb R\) is continuous and \(F\subseteq\mathbb R\) is closed. Since \(F\) is closed, its complement \( F^c\) is open. From part a, we have that \(f^{-1}(F^c)\) is open. Since \((f^{-1}(F))^c = f^{-1}(F^c)\), it follows that \( (f^{-1}(F))^c\) is open as well. Therefore \(f^{-1}(F)\) is closed.

    \item  Conversely, assume that \(f^{-1}(U)\) is open for every open set \(U \subseteq \mathbb{R}\). Let \(\epsilon > 0\) and \(c \in \mathbb{R}\). Consider the open set \(V_{\epsilon}(f(c))\). Since the preimage of every open set is open, it follows that \(f^{-1}(V_{\epsilon}(f(c)))\) is open. Moreover, since \(f(c) \in V_{\epsilon}(f(c))\) we have
          \[
              c \in f^{-1}(V_{\epsilon}(f(c))).
          \]
          Thus there exists \(\delta > 0\) such that
          \[
              V_{\delta}(c) \subseteq  f^{-1}(V_{\epsilon}(f(c)))
          \]

          Now suppose we have \(x \in \mathbb{R}\) such that \(x \in V_{\delta}(c)\).
          Then
          \[
              x \in f^{-1}(V_{\epsilon}(f(c))),
          \]
          so by definition of preimage,
          \[
              f(x) \in V_{\epsilon}(f(c)).
          \]
          That is, \(f\) is continuous as point \(c\).
          Since \( c \in \mathbb{R} \) was arbitrary, \(f\) is continuous on \(\mathbb{R}\).
\end{enumerate}

\end{document}